\section{Background and Related Work}
\label{sec:bg}

\subsection{The CGRA design space}

Coarse-grained reconfigurable arrays trade the bit-level flexibility of an FPGA
for word-level processing elements whose operation and interconnect are set by a
compact configuration~\cite{cgra-survey}. The design space is broad: PEs range
from single ALUs to tightly coupled VLIW datapaths, as in
ADRES~\cite{adres}; interconnects vary from nearest-neighbour meshes to richer
crossbars, as in MorphoSys~\cite{morphosys}; and configuration may be static for
a whole kernel or streamed as a sequence of contexts. More recent fabrics such
as Plasticine~\cite{plasticine} specialise the PEs and memories around parallel
patterns. Across this space the recurring tension is the same: richer PEs and
interconnect raise peak throughput but enlarge the configuration and the state
that must be managed each cycle.

The fabric described here sits at the minimal end of that spectrum by design. It
is a nearest-neighbour mesh of single-ALU, single-register PEs with a statically
loaded, one-word-per-PE configuration. This minimality is not incidental: it is
what lets the entire cycle-by-cycle behaviour of the array be exposed to, and
driven from, a host over an ordinary serial link, and it is what makes an exact
software model of the device tractable.

\subsection{Systolic dataflow}

The systolic principle---many identical cells, each doing a little work and
passing data to a neighbour on every beat, so that memory bandwidth is amortised
over a deep pipeline---dates to Kung~\cite{kung} and underpins modern matrix
engines such as the tensor processing unit~\cite{tpu}. A canonical
output-stationary array fuses, in each cell, a multiply with the accumulation of
a partial sum arriving from one neighbour while activations flow orthogonally,
so each cell holds two independent live values. Section~\ref{sec:dataflow} shows
why the single register per PE of this fabric precludes that fusion and selects
a \emph{weight-stationary} streaming dataflow instead---the same trade-off that
distinguishes weight- from output-stationary accelerators in the literature,
here forced by an explicit architectural constraint rather than chosen for
energy.

\subsection{Position of this work}

This report does not propose a new microarchitecture. Its contribution is a
complete, honest, end-to-end \emph{logical} treatment of a deliberately small
CGRA: a fully host-controlled fabric, a self-recovering byte-level protocol, a
byte-accurate device model that doubles as a golden reference, and a UNIX-style
software stack---all connected by a single narrow interface and all verifiable
without an FPGA. The value is pedagogical and methodological: it makes visible,
in one coherent system, how an architectural constraint propagates all the way
up to the dataflows a user can express.
